% =====================================================================
% PROPUESTA DE PROYECTO FINAL DE CARRERA
% Formato paper, dos columnas, lista para compilar con pdfLaTeX
% =====================================================================
\documentclass[9pt,twocolumn,a4paper]{extarticle}
% ------------------------- Idioma y tipografía -------------------------
\usepackage[utf8]{inputenc}
\usepackage[T1]{fontenc}
\usepackage[spanish,es-nodecimaldot]{babel}
\usepackage{times}
\usepackage{microtype}
% ------------------------- Página y columnas ---------------------------
\usepackage[
  a4paper,
  top=1.35cm,
  bottom=1.35cm,
  left=1.35cm,
  right=1.35cm,
  columnsep=0.55cm
]{geometry}
% ------------------------- Matemática y texto --------------------------
\usepackage{amsmath,amssymb}
\usepackage{enumitem}
\usepackage{titlesec}
\usepackage{graphicx}
\usepackage{xcolor}
\usepackage[hidelinks]{hyperref}
\usepackage[font=small,labelfont=bf,skip=2pt]{caption}
% ------------------------- TikZ ----------------------------------------
\usepackage{tikz}
\usetikzlibrary{
  arrows.meta,
  positioning,
  fit,
  calc,
  backgrounds
}
% ------------------------- Compactación --------------------------------
\setlist{nosep,leftmargin=1.1em}
\setlength{\parindent}{1em}
\setlength{\parskip}{0pt}
\setlength{\abovedisplayskip}{3pt}
\setlength{\belowdisplayskip}{3pt}
\setlength{\intextsep}{4pt}
\setlength{\textfloatsep}{5pt}
\setlength{\dbltextfloatsep}{5pt}
\setlength{\floatsep}{4pt}
\setcounter{topnumber}{4}
\setcounter{totalnumber}{6}
\setcounter{dbltopnumber}{4}
\renewcommand{\topfraction}{0.98}
\renewcommand{\textfraction}{0.02}
\renewcommand{\floatpagefraction}{0.90}
\renewcommand{\dbltopfraction}{0.98}
\renewcommand{\dblfloatpagefraction}{0.90}
\titlespacing*{\section}{0pt}{5pt}{2pt}
\titleformat{\section}{\normalfont\bfseries}{\thesection}{0.5em}{}
% =====================================================================
% PALETA DEL DIAGRAMA
% =====================================================================
\definecolor{ddBlue}{HTML}{2563EB}
\definecolor{ddCyan}{HTML}{0891B2}
\definecolor{ddGreen}{HTML}{3A8D24}
\definecolor{ddOrange}{HTML}{EA7A12}
\definecolor{ddPurple}{HTML}{8B5CF6}
\definecolor{ddRed}{HTML}{DC2626}
\definecolor{ddDark}{HTML}{172033}
\definecolor{ddGray}{HTML}{64748B}
\tikzset{
  module/.style={
    rectangle,
    rounded corners=7pt,
    fill=white,
    line width=1.1pt,
    minimum width=3.25cm,
    minimum height=6.25cm
  },
  groupblue/.style={
    rectangle,
    rounded corners=11pt,
    draw=ddBlue,
    fill=ddBlue!2,
    dash pattern=on 4pt off 3pt,
    line width=1pt
  },
  groupcyan/.style={
    rectangle,
    rounded corners=11pt,
    draw=ddCyan,
    fill=ddCyan!2,
    dash pattern=on 4pt off 3pt,
    line width=1pt
  },
  groupgreen/.style={
    rectangle,
    rounded corners=11pt,
    draw=ddGreen,
    fill=ddGreen!2,
    dash pattern=on 4pt off 3pt,
    line width=1pt
  },
  grouporange/.style={
    rectangle,
    rounded corners=11pt,
    draw=ddOrange,
    fill=ddOrange!2,
    dash pattern=on 4pt off 3pt,
    line width=1pt
  },
  header/.style={
    rectangle,
    rounded corners=4pt,
    text=white,
    font=\bfseries\footnotesize,
    inner xsep=9pt,
    inner ysep=5pt
  },
  cardtitle/.style={
    font=\bfseries\small,
    align=center,
    text width=2.85cm,
    text=ddDark
  },
  cardsubtitle/.style={
    font=\scriptsize,
    align=center,
    text width=2.75cm,
    text=ddDark
  },
  infobox/.style={
    rectangle,
    rounded corners=5pt,
    align=left,
    text width=2.72cm,
    inner sep=7pt,
    font=\tiny,
    line width=0.8pt
  },
  flow/.style={
    -{Latex[length=3mm,width=2mm]},
    draw=ddDark,
    line width=1.4pt
  }
}
% =====================================================================
% DOCUMENTO
% =====================================================================
\begin{document}
\pagestyle{empty}
\twocolumn[
\begin{center}
{\Large\bfseries
Representaciones auto-supervisadas para el descubrimiento de clases novedosas
en segmentación semántica de nubes de puntos 3D\par}
\vspace{4pt}
{\normalsize Renzo [Apellido] \quad Manyory Cueva Mendoza\par}
\vspace{2pt}
{\small Escuela Profesional de Ingeniería de Sistemas / Ciencia de la Computación\par}
\vspace{1pt}
{\small\itshape Propuesta de Proyecto Final de Carrera --- \today\par}
\end{center}
\vspace{3pt}
\begin{quotation}
\noindent\small\textbf{Resumen.}
La segmentación semántica de nubes de puntos opera bajo un supuesto de
conjunto cerrado: todo punto se asigna forzosamente a una clase conocida. El
paradigma \emph{Novel Class Discovery} (NCD) levanta ese supuesto, pero su
estado del arte en 3D ---HyperNCD~\cite{hyperncd}, 2026--- alcanza apenas
11.8--47.5 mIoU en clases novedosas frente a 49.8 del modelo supervisado en
SemanticKITTI. Identificamos una causa plausible: HyperNCD construye sus
prototipos con descriptores geométricos manuales (linealidad, planaridad y
\emph{scattering}). Sonata~\cite{sonata} demostró que apoyarse en señales
geométricas de bajo nivel induce el \emph{geometric shortcut}, un colapso
representacional que impide aprender semántica, y que corregirlo eleva el
\emph{linear probing} de 21.8\,\% a 72.5\,\% mIoU en ScanNet. Proponemos
sustituir el bloque geométrico manual de HyperNCD por el \emph{encoder}
auto-supervisado de Sonata, congelado, y medir su efecto en SemanticKITTI y
SemanticPOSS.
\vspace{2pt}
\noindent\small\textbf{Palabras clave:}
nubes de puntos, segmentación semántica 3D, \emph{novel class discovery},
aprendizaje auto-supervisado, hipergrafos.
\end{quotation}
\vspace{4pt}
]
% ---------------------------------------------------------------------
\section{Problema computacional y alcance}
% ---------------------------------------------------------------------
Partimos de nubes de puntos 3D ya capturadas y etiquetadas,
$P=\{(x_i,f_i)\}_{i=1}^{N}$, provistas por los conjuntos de datos
SemanticKITTI y SemanticPOSS. No se procesa señal cruda de sensor alguno,
sino únicamente los \emph{datasets} públicos. La segmentación semántica
asigna una etiqueta $y_i$ a cada punto. Los métodos actuales resuelven bien
la tarea ---PTv3~\cite{ptv3}: 77.5 mIoU; Sonata: 79.4 mIoU en ScanNet---,
pero comparten una restricción: \emph{el conjunto de clases está cerrado
durante el entrenamiento}. Un objeto no previsto no se marca como
``desconocido'', sino que se fuerza a la clase conocida más cercana, lo cual
constituye un modo de falla silencioso en conducción autónoma o robótica
móvil.

\emph{Novel Class Discovery} (NCD), introducido a nubes de puntos por Riz
\emph{et al.}~\cite{nops}, formaliza el problema: dado $\mathcal{C}_k$
(clases conocidas, etiquetadas) y $\mathcal{C}_n$ (clases novedosas, sin
etiqueta), el modelo debe segmentar ambas. El estado del arte,
HyperNCD~\cite{hyperncd}, modela asociaciones de alto orden mediante un
hipergrafo de \emph{prototipos geométrico-conscientes}. Cada punto se
representa como
$F_i=[f_{\mathrm{sem}}(x^i);f_{\mathrm{geo}}(x^i)]$, donde
$f_{\mathrm{geo}}$ se calcula analíticamente mediante linealidad,
planaridad y \emph{scattering} sobre $k$NN, con $k=15$.

Sonata~\cite{sonata} identifica el \emph{geometric shortcut}: el modelo
colapsa sobre pistas geométricas accesibles, como normales o altura, en lugar
de aprender semántica. Con \emph{linear probing}, PointContrast obtiene
5.6\,\%, MSC 21.8\,\% y Sonata 72.5\,\% mIoU en ScanNet. HyperNCD es posterior
a Sonata, pero emplea en su módulo central descriptores geométricos manuales.
A partir de esta relación se plantea la siguiente hipótesis:

\begin{quotation}
\noindent\footnotesize\itshape
Si $f_{\mathrm{geo}}$ se reemplaza por \emph{embeddings} de un
\emph{encoder} auto-supervisado congelado, los prototipos serán más
discriminativos y mejorará el mIoU de las clases novedosas.
\end{quotation}

\textbf{Incluye:} SemanticKITTI (secuencia 08), SemanticPOSS (secuencia 03),
cuatro \emph{splits} oficiales~\cite{nops,dasl}, \emph{encoder} congelado,
comparación contra NOPS, DASL e HyperNCD, y estudio de ablación.
\textbf{Excluye:} preentrenamiento SSL propio, detección de instancias,
vocabulario abierto basado en modelos de lenguaje y ejecución en tiempo real
sobre hardware embebido.

% ---------------------------------------------------------------------
\section{Diseño de proyecto}
% ---------------------------------------------------------------------
La Fig.~\ref{fig:diseno} presenta el flujo general del sistema:
\textbf{entrada $\rightarrow$ procesamiento $\rightarrow$ modelo
$\rightarrow$ salida}. El punto de intervención se encuentra en el módulo de
modelo: el \emph{encoder} congelado de Sonata reemplaza la rama geométrica
manual de HyperNCD antes de la concatenación de características. Todo lo
posterior ---formación de prototipos, construcción del hipergrafo dinámico y
generación de pseudo-etiquetas--- se mantiene para que la comparación sea
interpretable.

Los descriptores por punto se definen como
$f_{\mathrm{sem}}$ para la rama semántica de MinkowskiUNet-34C y
$f_{\mathrm{ssl}}$ para el \emph{encoder} auto-supervisado de Sonata. Ambos se
integran mediante
$F_i=[f_{\mathrm{sem}}(x^i);f_{\mathrm{ssl}}(x^i)]$.
La afinidad entre prototipos combina similitud geométrica y semántica:
$S_b=\alpha S_g+(1-\alpha)S_s$, donde $\alpha\in[0,1]$. Para cada prototipo
se seleccionan $M=8$ vecinos para formar las hiperaristas.


% =====================================================================
% DIAGRAMA DE DISEÑO DEL SISTEMA
% =====================================================================
\begin{figure*}[t]
\centering
\vspace{-3pt}
\resizebox{0.98\textwidth}{!}{%
\begin{tikzpicture}[font=\sffamily]

% Título del diagrama
\node[font=\bfseries\Large,text=ddDark] at (6.45,4.25)
{Diagrama de diseño del sistema propuesto};

% Tarjetas principales
\node[module,draw=ddBlue]   (entrada) at (0,0) {};
\node[module,draw=ddCyan]   (proceso) at (4.30,0) {};
\node[module,draw=ddGreen]  (modelo)  at (8.60,0) {};
\node[module,draw=ddOrange] (salida)  at (12.90,0) {};

% Contenedores externos
\begin{pgfonlayer}{background}
\node[groupblue,fit=(entrada),inner xsep=7pt,inner ysep=9pt] (gentrada) {};
\node[groupcyan,fit=(proceso),inner xsep=7pt,inner ysep=9pt] (gproceso) {};
\node[groupgreen,fit=(modelo),inner xsep=7pt,inner ysep=9pt] (gmodelo) {};
\node[grouporange,fit=(salida),inner xsep=7pt,inner ysep=9pt] (gsalida) {};
\end{pgfonlayer}

% Encabezados
\node[header,fill=ddBlue] at (gentrada.north) {MÓDULO 1: ENTRADA};
\node[header,fill=ddCyan] at (gproceso.north) {MÓDULO 2: PROCESAMIENTO};
\node[header,fill=ddGreen] at (gmodelo.north) {MÓDULO 3: MODELO};
\node[header,fill=ddOrange] at (gsalida.north) {MÓDULO 4: SALIDA};

% ---------------------------------------------------------------------
% MÓDULO 1: nube de puntos
% ---------------------------------------------------------------------
\begin{scope}[shift={($(entrada.north)+(0,-1.20)$)},scale=0.54]
\draw[ddBlue,line width=0.8pt]
(-1.35,-0.50)--(0,-1.25)--(1.35,-0.50)--(1.35,0.65)--
(0,1.40)--(-1.35,0.65)--cycle;
\draw[ddBlue!55,line width=0.6pt] (0,-1.25)--(0,0.05)--(1.35,0.65);
\draw[ddBlue!55,line width=0.6pt] (0,0.05)--(-1.35,0.65);
\draw[ddBlue!55,line width=0.6pt] (0,0.05)--(0,1.40);
\foreach \x/\y in {
-1.15/-0.30,-0.95/-0.42,-0.72/-0.55,-0.48/-0.70,-0.23/-0.86,
0.23/-0.86,0.48/-0.70,0.72/-0.55,0.95/-0.42,1.15/-0.30,
-1.02/0.08,-0.76/0.22,-0.50/0.36,-0.24/0.50,0/0.64,
0.24/0.50,0.50/0.36,0.76/0.22,1.02/0.08,
-0.70/0.68,-0.45/0.83,-0.20/0.98,0/1.10,0.20/0.98,0.45/0.83,0.70/0.68}
{\fill[ddBlue] (\x,\y) circle (1.6pt);}
\foreach \x/\h in {-0.75/0.70,-0.30/0.95,0.20/0.72,0.68/0.90}
{\draw[ddBlue,densely dotted,line width=1pt] (\x,-0.15)--(\x,\h);}
\end{scope}
\node[cardtitle] at ($(entrada.center)+(0,0.72)$) {Datos de entrada};
\node[cardsubtitle] at ($(entrada.center)+(0,0.08)$)
{Nubes de puntos 3D\\SemanticKITTI y SemanticPOSS};
\node[infobox,draw=ddBlue!65,fill=ddBlue!5]
at ($(entrada.center)+(0,-1.65)$)
{
\textbullet\ Coordenadas $(x,y,z)$\\[4pt]
\textbullet\ Intensidad o reflectancia\\[4pt]
\textbullet\ Clases conocidas etiquetadas\\[4pt]
\textbullet\ Clases novedosas sin etiqueta\\[4pt]
\textbullet\ Cuatro \emph{splits} oficiales
};

% ---------------------------------------------------------------------
% MÓDULO 2: preprocesamiento
% ---------------------------------------------------------------------
\begin{scope}[shift={($(proceso.north)+(0,-1.18)$)},scale=0.57]
\foreach \x/\y in {-0.70/1.10,0/1.10,0.70/1.10,-0.43/0.68,0.18/0.68,0.76/0.68}
{\fill[ddCyan] (\x,\y) circle (4pt);}
\filldraw[fill=ddCyan!80,draw=ddCyan,line width=0.8pt]
(-1.15,0.28)--(1.15,0.28)--(0.45,-0.48)--(0.45,-1.02)--
(0.14,-1.30)--(-0.14,-1.30)--(-0.45,-1.02)--(-0.45,-0.48)--cycle;
\fill[ddCyan] (-0.25,-1.62) circle (3.4pt);
\fill[ddCyan] (0.25,-1.62) circle (3.4pt);
\end{scope}
\node[cardtitle] at ($(proceso.center)+(0,0.72)$) {Preparación del \emph{input}};
\node[cardsubtitle] at ($(proceso.center)+(0,0.08)$)
{Transformación de la nube para\\el procesamiento computacional};
\node[infobox,draw=ddCyan!75,fill=ddCyan!5]
at ($(proceso.center)+(0,-1.65)$)
{
\textbullet\ Voxelización\\[4pt]
\textbullet\ Normalización\\[4pt]
\textbullet\ Muestreo y organización\\[4pt]
\textbullet\ MinkowskiEngine\\[4pt]
\textbullet\ Separación \emph{known}/\emph{novel}
};

% ---------------------------------------------------------------------
% MÓDULO 3: modelo
% ---------------------------------------------------------------------
\begin{scope}[shift={($(modelo.north)+(0,-1.16)$)},scale=0.55]
\foreach \ya in {0.85,0,-0.85}{
  \foreach \yb in {1.10,0.36,-0.36,-1.10}{
    \draw[ddGreen!65,line width=0.55pt] (-1.25,\ya)--(0,\yb);
  }
}
\foreach \yb in {1.10,0.36,-0.36,-1.10}{
  \draw[ddGreen!65,line width=0.55pt] (0,\yb)--(1.35,0);
}
\foreach \y in {0.85,0,-0.85}{
  \filldraw[fill=white,draw=ddGreen,line width=1.1pt]
  (-1.25,\y) circle (0.17);
}
\foreach \y in {1.10,0.36,-0.36,-1.10}{
  \filldraw[fill=white,draw=ddGreen,line width=1.1pt]
  (0,\y) circle (0.17);
}
\filldraw[fill=white,draw=ddGreen,line width=1.1pt]
(1.35,0) circle (0.17);
\end{scope}
\node[cardtitle] at ($(modelo.center)+(0,0.92)$)
{Segmentación y descubrimiento};
\node[cardsubtitle] at ($(modelo.center)+(0,0.16)$)
{Extracción de características y\\razonamiento entre clases};
\node[
  rectangle,
  rounded corners=3pt,
  draw=ddGreen,
  fill=ddGreen!10,
  text=ddGreen!55!black,
  font=\bfseries\tiny,
  inner xsep=5pt,
  inner ysep=3pt
] at ($(modelo.center)+(0,-0.32)$)
{APORTE: ENCODER SONATA};
\node[infobox,draw=ddGreen!75,fill=ddGreen!5]
at ($(modelo.center)+(0,-1.92)$)
{
\textbullet\ MinkowskiUNet-34C: $f_{\mathrm{sem}}$\\[4pt]
\textbullet\ Sonata/PTv3 congelado: $f_{\mathrm{ssl}}$\\[4pt]
\textbullet\ Integración: $F_i=[f_{\mathrm{sem}};f_{\mathrm{ssl}}]$\\[4pt]
\textbullet\ Prototipos por clase\\[4pt]
\textbullet\ Hipergrafo dinámico y pseudo-etiquetas
};

% ---------------------------------------------------------------------
% MÓDULO 4: salida segmentada
% ---------------------------------------------------------------------
\begin{scope}[shift={($(salida.north)+(0,-1.18)$)},scale=0.57]
\filldraw[fill=ddBlue!65,draw=black!70,line width=0.6pt]
(-1.35,-0.85) rectangle (-0.45,0.05);
\filldraw[fill=ddGreen!65,draw=black!70,line width=0.6pt]
(-0.40,-0.85) rectangle (0.45,0.20);
\filldraw[fill=ddOrange!75,draw=black!70,line width=0.6pt]
(0.50,-0.85) rectangle (1.35,0.15);
\filldraw[fill=ddPurple!65,draw=black!70,line width=0.6pt]
(-1.35,0.10) rectangle (-0.45,1.00);
\filldraw[fill=ddCyan!65,draw=black!70,line width=0.6pt]
(-0.40,0.25) rectangle (0.45,1.00);
\filldraw[fill=ddRed!65,draw=black!70,line width=0.6pt]
(0.50,0.20) rectangle (1.35,1.00);
\draw[black!80,line width=1pt] (-1.40,-0.90) rectangle (1.40,1.05);
\end{scope}
\node[cardtitle] at ($(salida.center)+(0,0.72)$) {Resultados del sistema};
\node[cardsubtitle] at ($(salida.center)+(0,0.08)$)
{Segmentación por punto de clases\\conocidas y novedosas};
\node[infobox,draw=ddOrange!80,fill=ddOrange!6]
at ($(salida.center)+(0,-1.65)$)
{
\textbullet\ Etiqueta semántica por punto\\[4pt]
\textbullet\ Clases conocidas y novedosas\\[4pt]
\textbullet\ Nube segmentada en 3D\\[4pt]
\textbullet\ mIoU \emph{novel}, \emph{known} y \emph{all}\\[4pt]
\textbullet\ Comparación con NOPS, DASL e HyperNCD
};

% Flechas y etiquetas
\draw[flow] (gentrada.east)--(gproceso.west);
\draw[flow] (gproceso.east)--(gmodelo.west);
\draw[flow] (gmodelo.east)--(gsalida.west);
\node[font=\tiny\itshape,fill=white,text=ddGray,inner sep=1pt]
at ($(gentrada.east)!0.5!(gproceso.west)+(0,0.24)$) {datos};
\node[font=\tiny\itshape,fill=white,text=ddGray,inner sep=1pt]
at ($(gproceso.east)!0.5!(gmodelo.west)+(0,0.24)$) {características};
\node[font=\tiny\itshape,fill=white,text=ddGray,inner sep=1pt]
at ($(gmodelo.east)!0.5!(gsalida.west)+(0,0.24)$) {predicciones};

\end{tikzpicture}%
}
\vspace{-2pt}
\caption{\textbf{Diagrama de diseño del sistema propuesto.}
Las nubes de puntos de SemanticKITTI y SemanticPOSS son voxelizadas y
normalizadas antes de ingresar al modelo. MinkowskiUNet-34C obtiene las
características semánticas $f_{\mathrm{sem}}$, mientras que el \emph{encoder}
auto-supervisado congelado de Sonata produce $f_{\mathrm{ssl}}$. Ambas
representaciones se integran para construir prototipos y un hipergrafo
dinámico que permite segmentar clases conocidas y descubrir clases novedosas.}
\label{fig:diseno}
\vspace{-7pt}
\end{figure*}

El criterio de éxito se fija sobre el \textbf{mIoU de clases novedosas} tras
emparejamiento húngaro: mejorar a HyperNCD en al menos dos de los cuatro
\emph{splits}, sin degradar el rendimiento de las clases conocidas.

\textbf{Riesgo y mitigación.}
Sonata se preentrenó mayoritariamente con interiores densos. En exteriores,
su \emph{linear probing} en SemanticKITTI puede quedar por debajo de un PTv3
entrenado desde cero. Como Plan B, se evaluará la concatenación
$F_i=[f_{\mathrm{sem}};f_{\mathrm{geo}};f_{\mathrm{ssl}}]$ en lugar del
reemplazo directo, permitiendo que el módulo de prototipos combine las tres
fuentes de información.

% ---------------------------------------------------------------------
\section{Objetivos del proyecto}
% ---------------------------------------------------------------------
\textbf{Objetivo general.}
Diseñar y evaluar un método de descubrimiento de clases novedosas para
segmentación 3D que sustituya los descriptores geométricos manuales por
representaciones auto-supervisadas preentrenadas, midiendo su impacto en
SemanticKITTI y SemanticPOSS.

\textbf{OE1: Reproducir el estado del arte.}
Replicar HyperNCD en ambos conjuntos de datos y en los cuatro \emph{splits}
oficiales, reportando la desviación sobre tres corridas.
\emph{Herramientas:} Python 3.10, PyTorch, MinkowskiEngine, código oficial de
HyperNCD y MinkowskiUNet-34C.
\emph{Alcance:} desviación máxima de $\pm1.0$ mIoU respecto de lo publicado.

\textbf{OE2: Integrar el \emph{encoder} SSL.}
Incorporar el \emph{encoder} de Sonata, con pesos congelados y
\emph{up-casting}, dentro del módulo de prototipos, en reemplazo del bloque
geométrico analítico.
\emph{Herramientas:} pesos públicos de Sonata, Pointcept, PTv3 y
\texttt{spconv}.
\emph{Alcance:} solo inferencia, sin retropropagación a través del
\emph{encoder}.

\textbf{OE3: Reajustar la configuración.}
Recalibrar $\alpha$ (balance geométrico-semántico), $M$ (vecinos por
hiperarista) y $\epsilon$ (suavizado de etiquetas).
\emph{Herramientas:} Optuna o búsqueda por grilla, AdamW y Weights\&Biases.
\emph{Alcance:} búsqueda en el \emph{split} 0 y aplicación sin reajuste a los
otros tres.

\textbf{OE4: Evaluar y realizar ablaciones.}
Comparar contra NOPS~\cite{nops}, DASL~\cite{dasl} e HyperNCD en mIoU
\emph{novel}, \emph{known} y \emph{all}; además, aislar la contribución del
\emph{encoder} mediante las variantes manual, auto-supervisada y concatenada.
\emph{Herramientas:} mIoU con emparejamiento húngaro y visualizaciones PCA.
\emph{Alcance:} cuatro \emph{splits} por conjunto de datos.

\textbf{Recursos.}
GPU con al menos 12 GB de VRAM. El \emph{encoder} congelado añade únicamente
costo de inferencia y permite almacenar en caché las características por
escena. Como respaldo se utilizará Colab Pro o Kaggle.

% ---------------------------------------------------------------------
\section{Resultados esperados}
% ---------------------------------------------------------------------
Se espera mejorar el mIoU de las clases novedosas respecto de HyperNCD en al
menos dos de los cuatro \emph{splits} de SemanticKITTI y SemanticPOSS, sin
degradar las clases conocidas. Los entregables serán: (i) tabla de
reproducción de HyperNCD con desviación reportada; (ii) implementación del
\emph{encoder} SSL dentro del módulo de prototipos; (iii) curvas de
sensibilidad de $\alpha$, $M$ y $\epsilon$; y (iv) tablas comparativas y un
estudio de ablación que aísle la contribución del \emph{encoder} congelado.

La contribución no consiste en crear una arquitectura completamente nueva,
sino en establecer empíricamente si el problema del \emph{geometric
shortcut}, demostrado en aprendizaje auto-supervisado 3D, también afecta el
descubrimiento de clases novedosas. Incluso un resultado nulo o negativo,
respaldado por ablaciones, constituirá un hallazgo reportable.

% ---------------------------------------------------------------------
% BIBLIOGRAFÍA
% ---------------------------------------------------------------------
{\small
\begin{thebibliography}{9}
\bibitem{nops}
L. Riz, C. Saltori, E. Ricci y F. Poiesi,
``Novel Class Discovery for 3D Point Cloud Semantic Segmentation'',
\emph{CVPR}, 2023.
\bibitem{ptv3}
X. Wu \emph{et al.},
``Point Transformer V3: Simpler, Faster, Stronger'',
\emph{CVPR}, 2024.
\bibitem{dasl}
R. Xu, C. Zhang, H. Ren y X. He,
``Dual-Level Adaptive Self-Labeling for Novel Class Discovery in Point Cloud
Segmentation'',
\emph{ECCV}, 2024.
\bibitem{sonata}
X. Wu \emph{et al.},
``Sonata: Self-Supervised Learning of Reliable Point Representations'',
\emph{arXiv:2503.16429}, 2025.
\bibitem{hyperncd}
Z. Zhang, A. Wu, Y. Li, Y. Han y J. Shen,
``Geometric-Aware Hypergraph Reasoning for Novel Class Discovery in Point
Cloud Segmentation'',
\emph{arXiv:2606.07280}, 2026.
\end{thebibliography}
}

\end{document}